\ExamPart{Trắc nghiệm nhiều phương án lựa chọn}{3.0}{Thí sinh trả lời từ câu 1 đến câu 12. Mỗi câu thí sinh chỉ chọn một phương án.}

\NQuestion Tam thức bậc hai $f(x)=2x^2-7x+3$ có hai nghiệm là
\ChoicesOneLine
{$x=\dfrac{1}{2}$ và $x=3$.}
{$x=1$ và $x=\dfrac{3}{2}$.}
{$x=-\dfrac{1}{2}$ và $x=3$.}
{$x=-1$ và $x=\dfrac{3}{2}$.}
{1}
\SolutionBlock{Ta có
\[
2x^2-7x+3=(2x-1)(x-3)=0.
\]
Suy ra hai nghiệm là $x=\dfrac{1}{2}$ và $x=3$.}

\NQuestion Tập nghiệm của bất phương trình $x^2-4x-5\le 0$ là
\ChoicesOneLine
{$[-1;5]$.}
{$(-\infty;-1]\cup[5;+\infty)$.}
{$(-1;5)$.}
{$(-\infty;-1)\cup(5;+\infty)$.}
{1}
\SolutionBlock{Ta có
\[
x^2-4x-5=(x+1)(x-5).
\]
Vì hệ số $a>0$ nên tam thức không dương trên đoạn giữa hai nghiệm. Vậy tập nghiệm là $[-1;5]$.}

\NQuestion Có bao nhiêu giá trị nguyên của $x$ thỏa mãn bất phương trình $(x-2)(x-6)<0$?
\ChoicesOneLine{$3$.}{$4$.}{$5$.}{$6$.}{1}
\SolutionBlock{Ta có $(x-2)(x-6)<0\iff 2<x<6$. Các giá trị nguyên là $3,4,5$, nên có $3$ giá trị.}

\NQuestion Một cửa hàng bán mỗi quyển sổ với giá $25\,000$ đồng. Chi phí sản xuất $x$ quyển sổ là $C(x)=x^2-20x+200$ (nghìn đồng), doanh thu là $R(x)=25x$ (nghìn đồng). Điều kiện để cửa hàng không lỗ là
\ChoicesTwoLines
{$5\le x\le 40$.}
{$x\le 5$ hoặc $x\ge 40$.}
{$5<x<40$.}
{$x<5$ hoặc $x>40$.}
{1}
\SolutionBlock{Không lỗ khi $R(x)\ge C(x)$, tức là
\[
25x\ge x^2-20x+200\iff x^2-45x+200\le 0.
\]
Ta có
\[
x^2-45x+200=(x-5)(x-40).
\]
Vì hệ số $a>0$, bất phương trình đúng khi
\[
5\le x\le 40.
\]
Vậy chọn đáp án A.}

\NQuestion Nghiệm của phương trình $\sqrt{x+1}=x-1$ là
\ChoicesOneLine{$x=2$.}{$x=3$.}{$x=2$ hoặc $x=3$.}{Vô nghiệm.}{2}
\SolutionBlock{Điều kiện: $x-1\ge0\iff x\ge1$.
Bình phương hai vế:
\[
x+1=(x-1)^2=x^2-2x+1\iff x^2-3x=0\iff x(x-3)=0.
\]
Đối chiếu điều kiện chỉ nhận $x=3$.}

\NQuestion Trong mặt phẳng tọa độ $Oxy$, cho $A(-1;2)$, $B(4;-3)$. Tọa độ vectơ $\overrightarrow{AB}$ là
\ChoicesOneLine{$(3;-5)$.}{$(5;-5)$.}{$(-5;5)$.}{$(5;1)$.}{2}
\SolutionBlock{Ta có
\[
\overrightarrow{AB}=(4-(-1);-3-2)=(5;-5).
\]}

\NQuestion Trung điểm của đoạn thẳng nối $M(2;5)$ và $N(-4;1)$ là
\ChoicesOneLine{$(-1;3)$.}{$(1;3)$.}{$(-2;6)$.}{$(3;2)$.}{1}
\SolutionBlock{Trung điểm là
\[
I\left(\dfrac{2+(-4)}{2};\dfrac{5+1}{2}\right)=(-1;3).
\]}

\NQuestion Đường thẳng $d:2x-y+3=0$ có một vectơ chỉ phương là
\ChoicesOneLine{$(2;-1)$.}{$(1;2)$.}{$(1;-2)$.}{$(2;1)$.}{2}
\SolutionBlock{Đường thẳng $2x-y+3=0$ có vectơ pháp tuyến $\vec n=(2;-1)$. Do đó một vectơ chỉ phương là $\vec u=(1;2)$.}

\NQuestion Phương trình đường thẳng đi qua $A(1;2)$ và có hệ số góc bằng $-3$ là
\ChoicesOneLine{$3x+y-5=0$.}{$3x-y+1=0$.}{$x+3y-7=0$.}{$3x+y+1=0$.}{1}
\SolutionBlock{Vì hệ số góc $m=-3$, nên
\[
y-2=-3(x-1)\iff 3x+y-5=0.
\]}

\NQuestion Phương trình đường tròn có tâm $I(2;-1)$ và bán kính $R=3$ là
\ChoicesTwoLines
{$(x-2)^2+(y+1)^2=9$.}
{$(x+2)^2+(y-1)^2=9$.}
{$(x-2)^2+(y-1)^2=3$.}
{$x^2+y^2-4x+2y+9=0$.}
{1}
\SolutionBlock{Đường tròn tâm $I(a;b)$ bán kính $R$ có phương trình
\[
(x-a)^2+(y-b)^2=R^2.
\]
Suy ra
\[
(x-2)^2+(y+1)^2=9.
\]}

\NQuestion Cho tam giác $ABC$ với $A(1;1)$, $B(5;1)$, $C(1;4)$. Tam giác $ABC$ là
\ChoicesOneLine{Tam giác vuông tại $A$.}{Tam giác vuông tại $B$.}{Tam giác cân tại $C$.}{Tam giác đều.}{1}
\SolutionBlock{Ta có $AB$ nằm ngang vì $y_A=y_B=1$, còn $AC$ thẳng đứng vì $x_A=x_C=1$. Do đó $AB\perp AC$, suy ra tam giác vuông tại $A$.}

\NQuestion Khoảng cách từ điểm $P(1;-2)$ đến đường thẳng $\Delta:x+2y-3=0$ bằng
\ChoicesOneLine{$\dfrac{6}{\sqrt5}$.}{$\dfrac{2}{\sqrt5}$.}{$\dfrac{5}{\sqrt5}$.}{$\sqrt5$.}{1}
\SolutionBlock{Ta có
\[
d(P,\Delta)=\dfrac{|1+2\cdot(-2)-3|}{\sqrt{1^2+2^2}}=\dfrac{6}{\sqrt5}.
\]}
